\documentclass[conference]{IEEEtran}

% =========================================================
% PLANTILLA LIGERA PARA OVERLEAF FREE
% Compilador recomendado: pdfLaTeX
% =========================================================

\usepackage[utf8]{inputenc}
\usepackage[T1]{fontenc}
\usepackage[spanish,es-nodecimaldot]{babel}
\usepackage{lmodern}
\usepackage{microtype}

\usepackage{amsmath,amssymb}
\usepackage{graphicx}
\usepackage{booktabs}
\usepackage{multirow}
\usepackage{array}
\usepackage{url}
\usepackage{cite}
\usepackage{xcolor}
\usepackage{listings}
\usepackage{algorithm}
\usepackage{algpseudocode}
\usepackage{hyperref}

\floatname{algorithm}{Algoritmo}

\definecolor{CodeBackground}{HTML}{F4F6F7}
\definecolor{CodeKeyword}{HTML}{005F73}
\definecolor{CodeComment}{HTML}{6C757D}
\definecolor{CodeString}{HTML}{BB3E03}

\hypersetup{
    colorlinks=true,
    linkcolor=black,
    citecolor=black,
    urlcolor=blue
}

\lstdefinestyle{paperstyle}{
    backgroundcolor=\color{CodeBackground},
    basicstyle=\ttfamily\footnotesize,
    keywordstyle=\color{CodeKeyword}\bfseries,
    commentstyle=\color{CodeComment}\itshape,
    stringstyle=\color{CodeString},
    numbers=left,
    numberstyle=\tiny,
    numbersep=5pt,
    frame=single,
    breaklines=true,
    showstringspaces=false,
    tabsize=4
}

\lstset{
    style=paperstyle,
    language=Java
}

% =========================================================
% DATOS DEL ARTÍCULO
% =========================================================

\title{
Sistema de Priorización de Atención: Diseño e Implementación Inicial de una
Cola de Prioridad por Niveles y una Pila sobre Listas Enlazadas
}

\author{
\IEEEauthorblockN{
Johan Adrian Torres Robles,
Esteban Mora Lopez,
Miguel Angel Rodriguez Bohorquez
}
\IEEEauthorblockA{
\textit{Universidad Nacional de Colombia}\\
\textit{Sede Bogotá}\\
Bogotá D.C., Colombia\\
\{jotorresro, emoral, mrodriguezbo\}@unal.edu.co
}
}
\begin{document}

\maketitle

% =========================================================
% RESUMEN
% =========================================================

\begin{abstract}
Los centros de atención que reciben solicitudes de servicio deben atenderlas
según su nivel de urgencia y no según su orden de llegada, además de poder
consultar y cancelar solicitudes puntuales y revisar las atenciones
realizadas. Este documento presenta el diseño e implementación inicial de un
Sistema de Priorización de Atención (SPA) que resuelve este problema con
estructuras lineales. Una cola de prioridad formada por una cola FIFO por
nivel de urgencia, cada una sobre una lista enlazada con referencias a la
cabeza y a la cola, registra y atiende solicitudes en tiempo constante y
resuelve el desempate por orden de llegada sin comparaciones adicionales; la
búsqueda y la cancelación por identificador recorren estas listas en tiempo
lineal; y una pila enlazada registra el historial de atenciones. La primera
entrega incluye una implementación funcional con menú de consola que cubre
el registro, la atención, la búsqueda, la cancelación y la consulta del
historial, verificada mediante un programa de prueba manual con casos
borde. Deshacer la última atención y la reducción del costo de la búsqueda
quedan planteados para las siguientes entregas.
\end{abstract}

\begin{IEEEkeywords}
cola de prioridad, cola, lista enlazada, pila, estructuras de datos
\end{IEEEkeywords}

% =========================================================
% CONTENIDO
% =========================================================

\section{Introducción}

Los centros de atención que reciben solicitudes de servicio (salas de
urgencias, mesas de ayuda técnica, módulos de atención al cliente) enfrentan
un mismo reto operativo: las solicitudes no deben atenderse en el orden en
que llegan, sino según su nivel de urgencia, y el personal necesita poder
consultar o cancelar solicitudes específicas y revisar las atenciones
realizadas sin tener que reorganizar manualmente toda la fila.

Resolver este problema de manera eficiente requiere estructuras de datos que
soporten, simultáneamente: (i) obtener en todo momento la solicitud más
urgente sin tener que ordenar repetidamente toda la colección, (ii) respetar
el orden de llegada entre solicitudes de igual urgencia, (iii) ubicar y
retirar una solicitud puntual por su identificador, y (iv) registrar las
atenciones en el orden en que ocurren para poder consultarlas y, en una
entrega posterior, revertir la más reciente.

Este documento presenta el diseño e implementación inicial de un
\emph{Sistema de Priorización de Atención} (SPA), que resuelve el problema
anterior combinando estructuras lineales estudiadas en la asignatura: colas
y listas enlazadas para el orden de atención, y una pila para el historial
de atenciones.

La contribución de esta primera entrega consiste en: (a) la formulación del
problema y su alcance, (b) la selección y justificación técnica de las
estructuras de datos empleadas, y (c) una implementación funcional que
permite registrar, atender, buscar y cancelar solicitudes, y consultar el
historial de atenciones.

El resto del documento se organiza así: la Sección~\ref{sec:relacionado}
presenta trabajo relacionado; la Sección~\ref{sec:metodologia} describe los
objetivos, el alcance, los requerimientos y el diseño de la solución,
incluyendo las estructuras de datos seleccionadas; la
Sección~\ref{sec:resultados} reporta el estado actual de la implementación;
la Sección~\ref{sec:discusion} discute las limitaciones conocidas; y la
Sección~\ref{sec:conclusiones} resume el estado del proyecto y el trabajo
futuro.

\section{Trabajo relacionado}
\label{sec:relacionado}

El problema de atender elementos según su prioridad se modela con el tipo
abstracto \emph{cola de prioridad}, cuyas operaciones básicas son insertar
un elemento y extraer el de mayor prioridad~\cite{cormen,goodrich}. La
literatura distingue entre el tipo abstracto y su representación: el mismo
comportamiento puede implementarse con estructuras distintas, y es la
representación la que determina el costo de cada operación.

Una representación directa es una lista ordenada por prioridad, que permite
extraer el elemento más prioritario en tiempo constante, pero cuya inserción
cuesta $O(n)$, pues debe recorrer la lista hasta encontrar la posición
correcta. En el extremo opuesto, una lista sin orden inserta en tiempo
constante pero obliga a recorrerla completa, en $O(n)$, cada vez que se
extrae el elemento más prioritario~\cite{goodrich}. Para prioridades
arbitrarias, la representación clásica que equilibra ambos costos es el
\emph{heap} binario, con inserción y extracción en $O(\log n)$~\cite{cormen}.

Cuando la prioridad solo puede tomar un número pequeño y fijo de valores,
como ocurre en los esquemas de triage que clasifican a cada paciente en
unos pocos niveles de urgencia, existe una alternativa más simple: mantener
una cola FIFO independiente por cada nivel. Insertar consiste en encolar en
el nivel correspondiente, y extraer consiste en revisar los niveles en orden
de urgencia hasta encontrar uno no vacío. Como el número de niveles es una
constante, ambas operaciones cuestan $O(1)$, y el orden FIFO de cada cola
resuelve de forma natural el desempate por orden de llegada. Esta es la
representación que adopta este proyecto (Sección~\ref{sec:metodologia}).

Sistemas reales de mesa de ayuda y de triage hospitalario resuelven un
problema equivalente: priorizar solicitudes por urgencia y permitir su
consulta puntual. Este proyecto no pretende replicar un sistema de
producción, sino aislar y resolver correctamente el núcleo algorítmico
común a este tipo de sistemas.

\section{Metodología}
\label{sec:metodologia}

\subsection{Objetivos}

\textbf{Objetivo general.} Diseñar e implementar un sistema que permita
registrar, priorizar, atender, cancelar y consultar solicitudes de atención
de forma eficiente, aplicando estructuras de datos apropiadas para cada
operación.

\textbf{Objetivos específicos.}
\begin{enumerate}
  \item Implementar una cola de prioridad formada por una cola FIFO por
    nivel de urgencia, cada una sobre una lista enlazada, para determinar
    la solicitud a atender en tiempo constante.
  \item Implementar la búsqueda y la cancelación de una solicitud pendiente
    por su identificador.
  \item Implementar una pila sobre lista enlazada que registre el historial
    de atenciones y permita, en una entrega posterior, deshacer la más
    reciente.
  \item Integrar las estructuras en una interfaz de consola que exponga las
    operaciones del sistema.
\end{enumerate}

\subsection{Alcance y limitaciones}

El sistema completo contempla: registrar una solicitud con un
identificador, un nombre y un nivel de urgencia (de 1, más urgente, a 5,
menos urgente); atender la solicitud de mayor prioridad; buscar y cancelar
una solicitud pendiente por identificador; consultar el historial de
atenciones; y deshacer la atención más reciente. Esta primera entrega
implementa todas estas operaciones excepto deshacer, que queda diseñada y
se implementará en la siguiente entrega.

Quedan explícitamente fuera del alcance del proyecto: persistencia en base
de datos, interfaz gráfica, acceso concurrente de múltiples usuarios y
autenticación. Se asume que los identificadores son proporcionados por
quien registra la solicitud; el sistema rechaza un identificador que ya
pertenezca a otra solicitud pendiente. Se asume también que la urgencia
toma solo cinco valores, supuesto del que depende el costo constante de la
cola de prioridad elegida.

\subsection{Requerimientos principales}

\begin{itemize}
  \item RF1: El sistema debe permitir registrar una solicitud con
    identificador único entre las pendientes, nombre y nivel de urgencia
    entre 1 y 5, rechazando valores fuera de ese rango.
  \item RF2: El sistema debe atender siempre la solicitud pendiente de
    mayor urgencia; en caso de empate, debe atenderse primero la que llegó
    antes.
  \item RF3: El sistema debe permitir buscar una solicitud pendiente por su
    identificador.
  \item RF4: El sistema debe permitir cancelar una solicitud pendiente de
    modo que nunca sea atendida.
  \item RF5: El sistema debe permitir deshacer la atención más reciente,
    devolviendo la solicitud correspondiente a la cola de pendientes.
  \item RF6: El sistema debe permitir consultar el historial de
    solicitudes atendidas.
\end{itemize}

\subsection{Diseño de la solución}

El sistema se organiza en los siguientes componentes: una clase
\texttt{Ticket} que representa una solicitud; una cola genérica
\texttt{LinkedQueue<T>} sobre lista enlazada; una cola de prioridad
\texttt{LevelPriorityQueue}, formada por un arreglo de cinco
\texttt{LinkedQueue<Ticket>}, una por nivel de urgencia; una pila genérica
\texttt{LinkedStack<T>}; y un historial \texttt{ServiceHistory} construido
sobre esa pila. Una clase fachada, \texttt{TriageSystem}, coordina la cola
de prioridad y el historial para que cada operación pública los mantenga
consistentes, y una clase \texttt{Main} ofrece el menú de consola.

\begin{figure}[h]
  \centering
  \fbox{\parbox{0.9\linewidth}{\centering
    \texttt{Main} (consola) \\[4pt]
    $\downarrow$ \\[4pt]
    \texttt{TriageSystem} \\[4pt]
    $\swarrow$ \hspace{3em} $\searrow$ \\[4pt]
    \texttt{LevelPriorityQueue} \hspace{1em} \texttt{ServiceHistory} \\[2pt]
    {\footnotesize 5 $\times$ \texttt{LinkedQueue<Ticket>}} \hspace{2.5em}
    {\footnotesize \texttt{LinkedStack<Ticket>}}
  }}
  \caption{Componentes principales del sistema y la estructura de datos
  que resuelve cada responsabilidad.}
  \label{fig:arquitectura}
\end{figure}

\subsection{Estructuras de datos seleccionadas}

\textbf{Cola de prioridad por niveles.} Representa el conjunto de
solicitudes pendientes. Está formada por un arreglo de cinco colas FIFO;
la cola de índice $u-1$ contiene, en orden de llegada, las solicitudes de
urgencia $u$. Debe soportar insertar una solicitud y extraer la más
prioritaria. Insertar consiste en encolar al final de la cola de su nivel,
y extraer consiste en recorrer los niveles desde el más urgente y desencolar
del primero que no esté vacío (Algoritmo~\ref{alg:atender}). Como el número
de niveles es la constante $k=5$, ambas operaciones cuestan $O(k)=O(1)$. El
desempate por orden de llegada (RF2) no requiere comparaciones adicionales:
lo garantiza el comportamiento FIFO de cada cola. Se descartó una lista
ordenada por prioridad, cuya inserción cuesta $O(n)$, y no se requiere un
\emph{heap}, cuyo costo $O(\log n)$ solo se justifica cuando la prioridad
puede tomar muchos valores distintos.

\textbf{Cola sobre lista enlazada con referencias a cabeza y cola.} Cada
nivel es una cola implementada como lista simplemente enlazada que mantiene
una referencia a la cabeza (frente) y otra al último nodo (final). La
referencia al final es la que permite encolar en $O(1)$; sin ella, encolar
requeriría recorrer la lista en $O(n)$. Se prefirió la lista enlazada sobre
un arreglo circular porque el número de solicitudes por nivel no se conoce
de antemano y porque la cancelación necesita retirar un elemento
arbitrario, operación que en una lista enlazada solo exige reenlazar el
nodo anterior con el siguiente, sin desplazar elementos.

\textbf{Búsqueda y cancelación por identificador.} La cola de prioridad
resuelve eficientemente ``\textquestiondown{}cuál es la siguiente solicitud a
atender?'', pero no ``\textquestiondown{}dónde está la solicitud con
identificador X?''. En esta entrega la búsqueda recorre las colas de los
cinco niveles, con costo $O(n)$ en el peor caso. La cancelación usa esa
búsqueda y luego elimina el nodo de la lista enlazada de su nivel en $O(1)$,
actualizando la referencia al final cuando el nodo eliminado es el último
(Algoritmo~\ref{alg:cancelar}); su costo total es $O(n)$. La misma búsqueda
permite rechazar identificadores repetidos al registrar, por lo que el
registro completo cuesta $O(n)$. Reducir estos costos a $O(1)$ esperado
mediante un índice por identificador es una de las mejoras previstas para
las siguientes entregas (Sección~\ref{sec:discusion}).

\textbf{Pila (lista enlazada simple, encabezado único).} Registra las
solicitudes ya atendidas en orden inverso de atención. Se eligió una pila
porque la operación prevista, deshacer la atención más reciente (RF5), es
exactamente el comportamiento LIFO: apilar al atender, en $O(1)$, y
desapilar al deshacer, en $O(1)$. Se implementó como lista enlazada con
solo referencia al tope, ya que todas las operaciones ocurren en ese
extremo. Consultar el historial completo (RF6) recorre la pila en $O(h)$,
siendo $h$ el número de atenciones.

En los algoritmos, $L[1..5]$ denota el arreglo de colas por nivel y $S$ la
pila del historial.

\begin{algorithm}[t]
\caption{Atender la siguiente solicitud}
\label{alg:atender}
\begin{algorithmic}[1]
\Function{ServeNext}{}
  \For{$u \gets 1$ \textbf{to} $5$}
    \If{$L[u]$ no está vacía}
      \State $t \gets$ \Call{Dequeue}{$L[u]$}
      \State apilar $t$ en $S$
      \State \Return $t$
    \EndIf
  \EndFor
  \State \Return \textsc{null}
\EndFunction
\end{algorithmic}
\end{algorithm}

\begin{algorithm}[t]
\caption{Cancelar una solicitud pendiente}
\label{alg:cancelar}
\begin{algorithmic}[1]
\Function{CancelRequest}{$id$}
  \For{$u \gets 1$ \textbf{to} $5$}
    \State $prev \gets$ \textsc{null}; $cur \gets L[u].head$
    \While{$cur \neq$ \textsc{null}}
      \If{$cur.value.id = id$}
        \If{$prev =$ \textsc{null}}
          \State $L[u].head \gets cur.next$
        \Else
          \State $prev.next \gets cur.next$
        \EndIf
        \If{$cur = L[u].tail$}
          \State $L[u].tail \gets prev$
        \EndIf
        \State \Return \textsc{true}
      \EndIf
      \State $prev \gets cur$; $cur \gets cur.next$
    \EndWhile
  \EndFor
  \State \Return \textsc{false}
\EndFunction
\end{algorithmic}
\end{algorithm}

\section{Estado de la implementación}
\label{sec:resultados}

La Tabla~\ref{tab:estado} resume el estado de cada funcionalidad a la
fecha de esta entrega, distinguiendo entre funcionalidades implementadas,
parcialmente implementadas y pendientes.

\begin{table}[t]
  \caption{Estado de las funcionalidades en la primera entrega}
  \label{tab:estado}
  \centering
  \begin{tabular}{p{0.58\columnwidth}p{0.28\columnwidth}}
    \toprule
    Funcionalidad & Estado \\
    \midrule
    Registrar solicitud (RF1) & Implementada \\
    Atender por prioridad (RF2) & Implementada \\
    Buscar por identificador (RF3) & Implementada \\
    Cancelar solicitud (RF4) & Implementada \\
    Deshacer última atención (RF5) & Pendiente \\
    Consultar historial (RF6) & Implementada \\
    Interfaz de consola & Implementada \\
    Validación de entradas & Parcial \\
    Pruebas del sistema & Parcial \\
    Índice por identificador en $O(1)$ & Pendiente \\
    Análisis experimental de desempeño & Pendiente \\
    \bottomrule
  \end{tabular}
\end{table}

\textbf{Implementadas.} Cinco de los seis requerimientos funcionales están
disponibles en la clase \texttt{TriageSystem}: \texttt{registerRequest},
\texttt{serveNext}, \texttt{findById}, \texttt{cancelRequest} y
\texttt{getHistory}. La interfaz de consola (\texttt{Main}) expone todas
estas operaciones mediante un menú. Las estructuras de soporte
(\texttt{LinkedQueue}, \texttt{LevelPriorityQueue} y \texttt{LinkedStack})
son implementación propia; de la biblioteca estándar de Java solo se usa
\texttt{ArrayList} para devolver el historial como lista.

\textbf{Parcialmente implementadas.} La validación de entradas rechaza
identificadores vacíos o repetidos entre las solicitudes pendientes,
urgencias fuera del rango 1--5 y urgencias no numéricas, pero aún no valida
el nombre de la solicitud. Las pruebas se realizan mediante un programa de
prueba manual (\texttt{BasicTests.java}) que imprime el resultado esperado
junto al obtenido, sin aserciones automáticas. El escenario principal
registra cuatro solicitudes, verifica el orden de atención por urgencia y
el desempate por orden de llegada entre dos solicitudes de urgencia 1,
cancela una solicitud pendiente y comprueba que deja de encontrarse y de
contarse, y confirma el orden del historial. Los casos borde cubren una
urgencia fuera de rango, un identificador repetido, la cancelación de un
identificador inexistente, la reutilización del identificador de una
solicitud ya atendida y la atención con la cola vacía.

\textbf{Pendientes.} Para la siguiente entrega quedan deshacer la última
atención (RF5), cuyo diseño se apoya en la pila ya implementada; un índice
por identificador que reduzca la búsqueda y la cancelación a tiempo
esperado constante; y la migración de las pruebas manuales a pruebas
automatizadas. Para la entrega final queda un análisis experimental de
tiempos de ejecución sobre volúmenes de datos mayores.

No se reportan resultados experimentales de desempeño en esta entrega,
dado que el énfasis de la primera versión es la corrección funcional y la
justificación del diseño; el análisis de complejidad teórica se presentó
en la Sección~\ref{sec:metodologia}.

\section{Discusión}
\label{sec:discusion}

El diseño adoptado separa la responsabilidad de cada estructura: la cola de
prioridad por niveles decide el orden de atención, la lista enlazada de
cada nivel permite retirar solicitudes canceladas, y la pila registra el
historial. Esta separación evita resolver todas las necesidades con una
sola estructura, por ejemplo una lista que se recorra linealmente para todo,
lo que degradaría también la atención a $O(n)$.

La principal limitación conocida es el costo lineal de la búsqueda por
identificador, que se traslada a la cancelación y a la verificación de
identificadores repetidos al registrar. Para el tamaño de datos esperado
(una fila de atención, no miles de solicitudes simultáneas) este costo se
considera aceptable en esta etapa. Una mejora prevista es mantener un
índice por identificador que permita ubicar cada solicitud en tiempo
esperado constante sin alterar el orden de atención.

Otra limitación es que el costo constante de la cola de prioridad depende
de que la urgencia tome solo cinco valores. Si en el futuro la prioridad
dependiera de una magnitud continua, como el tiempo de espera acumulado,
habría que evaluar otra representación de la cola de prioridad.

Finalmente, los datos existen solo durante la ejecución del programa, lo
cual es consistente con el alcance definido
(Sección~\ref{sec:metodologia}).

\section{Conclusiones}
\label{sec:conclusiones}

Esta primera entrega formuló el problema de priorizar y gestionar
solicitudes de atención y presentó una solución basada en estructuras
lineales: una cola de prioridad formada por una cola enlazada por nivel de
urgencia, que registra y atiende en tiempo constante y resuelve el
desempate por orden de llegada sin comparaciones adicionales, y una pila
enlazada para el historial de atenciones. La implementación inicial cubre
cinco de los seis requerimientos funcionales y fue verificada mediante un
programa de prueba manual que recorre el escenario principal y varios casos
borde.

Como trabajo futuro se plantea: implementar el deshacer de la última
atención sobre la pila existente, incorporar un índice por identificador
para reducir el costo de la búsqueda y la cancelación, añadir pruebas
automatizadas y medir experimentalmente el desempeño con volúmenes de datos
mayores.


% =========================================================
% AGRADECIMIENTOS
% =========================================================

\section*{Agradecimientos}

Los autores agradecen a la Universidad Nacional de Colombia y al equipo
docente de la asignatura Estructuras de Datos.

% =========================================================
% REFERENCIAS
% Sin BibTeX ni Biber para reducir tiempo de compilación.
% =========================================================

\begin{thebibliography}{00}

\bibitem{cormen}
T. H. Cormen, C. E. Leiserson, R. L. Rivest y C. Stein,
\textit{Introduction to Algorithms}, 4th ed.
Cambridge, MA, USA: MIT Press, 2022.

\bibitem{goodrich}
M. T. Goodrich, R. Tamassia y M. H. Goldwasser,
\textit{Data Structures and Algorithms in Java}, 6th ed.
Hoboken, NJ, USA: Wiley, 2014.

\end{thebibliography}

\end{document}
